\documentclass[12pt]{article}

\input{../../_assets/preambulo.tex}


\begin{document}

    % 1. Foto de fondo
    % 2. Título
    % 3. Encabezado Izquierdo
    % 4. Color de fondo
    % 5. Coord x del titulo
    % 6. Coord y del titulo
    % 7. Fecha

    
    \input{../../_assets/portada}
    \portadaExamen{ffccA4.jpg}{Ecuaciones\\Diferenciales II\\Examen XXIV}{Ecuaciones Diferenciales II. Examen XXIV}{MidnightBlue}{-8}{28}{2026}{}

    \begin{description}
        \item[Asignatura] Ecuaciones Diferenciales II.
        \item[Curso Académico] 2024/2025.
        \item[Grado] Doble Grado en Ingeniería Informática y Matemáticas.
        \item[Grupo] Único.
        \item[Profesor] Rafael Ortega Ríos.
        \item[Descripción] Convocatoria Ordinaria.
        \item[Fecha] 12 de junio de 2025.
        % \item[Duración] ---.
    \end{description}
    \newpage


    % ------------------------------------
    
    \begin{ejercicio}
        Se considera la sucesión de funciones
        \[
            f_n : \left[0, 1\right] \longrightarrow \bb{R}, \quad f_n(t) = t^n(1-t).
        \]
        \begin{enumerate}[label=\alph*)]
            \item Calcula $\|f_n\|_\infty = \max\limits_{t \in \left[0, 1\right]} |f_n(t)|$.
            \item ¿Es convergente la sucesión $\{f_n\}$? ¿Es equicontinua?
        \end{enumerate}
    \end{ejercicio}

    \begin{ejercicio}
        Prueba que existen dos funciones continuas $x, y : \bb{R} \longrightarrow \bb{R}$, $x = x(t)$, $y = y(t)$ tales que se cumple, para cada $t \in \bb{R}$,
        \[
            x(t) = 7 + \int_{2}^{t} \frac{y(s)}{1 + y(s)^2} \, ds, \quad y(t) = 8 + \int_{2}^{t} e^s x(s) \, ds.
        \]
        ¿Son únicas estas funciones?
    \end{ejercicio}

    \begin{ejercicio}
        Para cada $n \geq 1$ se considera la solución maximal $x_n(t)$ del problema de valores iniciales
        \[
            \dot{x} = x^5, \quad x(0) = 1 + \frac{1}{n}.
        \]
        Calcula, si existe, $\lim\limits_{n \to \infty} \int_{0}^{1} x_n(s)^2 \, ds$.
    \end{ejercicio}

    \begin{ejercicio}
        Se considera el problema de valores iniciales
        \[
            \ddot{x} - y^3 = 0, \quad \dot{y} = \cos x, \quad x(0) = 0, \quad \dot{x}(0) = 1, \quad y(0) = 1.
        \]
        \begin{enumerate}[label=\alph*)]
            \item Demuestra que existe una única solución maximal.
            \item ¿Está definida esta solución en toda la recta real?
        \end{enumerate}
    \end{ejercicio}

    \begin{ejercicio}
        Para cada $n \geq 10^6$ se supone que $x_n : \left]-1, \nicefrac{1}{2}\right[ \longrightarrow \bb{R}$ es una función continua que cumple
        \[
            x_n(t) = \frac{1}{n} + \int_{0}^{t} e^{x_n(s)} \, ds, \quad t \in \left]-1, \nicefrac{1}{2}\right[.
        \]
        Demuestra que la sucesión $\{x_n\}$ es uniformemente convergente y calcula su límite.
    \end{ejercicio}

    \begin{ejercicio}
        Dado $\veps \in \bb{R}$ se considera el problema de valores iniciales
        \[
            \ddot{x} = -x - x^3 y^4, \quad \ddot{y} = -y - x^4 y^3, \quad x(\veps) = \veps, \quad y(\veps) = 0, \quad \dot{x}(\veps) = 0, \quad \dot{y}(\veps) = \veps.
        \]
        \begin{enumerate}[label=\alph*)]
            \item Prueba que existe una única solución maximal $(x(t, \veps), y(t, \veps))$ definida para $t \in \left]-\infty, +\infty\right[$.
            \item Encuentra dos funciones $a(t)$ y $b(t)$ tales que, para cada $t \in \bb{R}$,
            \[
                y(t, \veps) = a(t) + b(t)\veps + o(\veps), \quad \text{si } \veps \to 0.
            \]
        \end{enumerate}
    \end{ejercicio}

    \newpage
    % \setcounter{ejercicio}{0} % Reiniciar contador de ejercicios
    % \noindent
    % \textbf{Solución.}

\end{document}